\scanpage{31}
\mspage{25}

\textcircled{a}
\[
 \deleted{(W^{\mathrm{univ}})_*\subset\mathrm{Fib}_W}\ \text{et}\quad
 (W^{\mathrm{univ}})_*\cap\unsure{\overline{W}|W}^{\deleted{\unsure{\mathrm{univ}}}}_{\mathrm{univ}}=(\mathrm{Cof}_{\unsure{W|M}})_*
\]

\textcircled{b}
\[
 \deleted{\left\{\begin{aligned}
 (W^{\mathrm{univ}})_*&\subset\mathrm{Fib}_W\ \text{et}\\
 (\mathrm{Cof}_W)_*&\subset W_{\mathrm{univ}}
 \end{aligned}\right.}\qquad
 \underbrace{\widetilde{W^{\mathrm{univ}}\cup(W_{\mathrm{univ}})^*}=\mathrm{Cof}_W}
 _{\substack{\text{i.e. }\mathrm{Cof}_W\text{ est }\deleted{\text{\unsure{engendré}}}\inserted{\text{la partie}}\\
 \text{close à g. engendrée par}\\W^{\mathrm{univ}}\ \text{et}\ (W_{\mathrm{univ}})^*}},
\]

et ces conditions impliquent \inserted{\textcircled{c}} qu'il existe
une structure de \uline{triple \deleted{\illegible{1}}\inserted{clos} de Quillen}
$(\deleted{\unsure{F}}W,C,F)$ induisant $W$ (donc avec
$TC=W\cap C$, $TF=W\cap F$) (lequel sera clos au
sens de Quillen ssi $W$ stable par
facteurs directs --- car $C$ et $F$ sont déjà
stables spontanément).

Il suffit en effet de prendre donc
\[
 TC=W^{\mathrm{univ}},\quad C=\mathrm{Cof}_W.
\]
Mais je doute que \unsure{-} l'existence d'un
triple clos de Quillen $(W,C,F)$ suffise
à impliquer les conditions a), b)\unsure{$\ne$|,|?}
ci-dessus, ni $=$ la condition plus faible
\[
 (W^{\mathrm{univ}})_*\cap W_{\mathrm{univ}}=(\mathrm{Cof}_W)_*
\]
qui \deleted{suffit déjà} signifie \uline{exactement} que
$C=\mathrm{Cof}_W$, $TC=W_{\deleted{\unsure{\mathrm{univ}}}}^{\inserted{\mathrm{univ}}}$ et leurs transposées
$TF=(\mathrm{Cof}_W)_*$, $F=(W^{\mathrm{univ}})_*$ forment une \deleted{\unsure{syst.}} catégorie
de Quillen close
% archive sheet number: 30
